% Standalone research draft. FL-to-FU tension and deletion-boundary revision, 2026-10-14.
% R028--R031 are saved-model audits, not newly trained independent seeds.
\documentclass[10pt,a4paper]{article}
\usepackage{kotex}
\usepackage{fontspec}
\IfFontExistsTF{Noto Sans CJK KR}{\setmainhangulfont{Noto Sans CJK KR}}{\setmainhangulfont{Malgun Gothic}}
\usepackage[margin=23mm]{geometry}
\usepackage{amsmath,amssymb,amsthm,booktabs,tabularx,array}
\usepackage[colorlinks=true,linkcolor=blue,citecolor=blue,urlcolor=blue]{hyperref}
\newtheorem{proposition}{명제}
\newcommand{\ind}{\mathbf1}
\newcommand{\pos}[1]{\left[#1\right]_+}
\newcommand{\med}{\operatorname{median}}
\title{FL에서 잠복하고 FU에서 활성화되는 백도어의 탐지 연구:\\업데이트 이력과 삭제 요청에 따른 상쇄 해제}
\author{연구 초안}
\date{v7 개정 / 탐지 우선 및 실험 근거 감사 / 2026년 10월 15일}
\begin{document}
\maketitle

\begin{abstract}
현재의 주 검증 대상은 삭제 요청에 따른 상쇄 해제 점수가 단순 선형 삭제 영향보다
독립 정상 non-IID 조건에서 추가 탐지 정보를 제공하는가이다.
이하 제어 설계는 후속 연구이며, 제어 성공이나 불수락을 탐지 성공으로 세지 않는다.
본 연구의 고정된 대상은 FL에서 심은 백도어가 FU를 통해 활성화되는 공격이다.
장력은 그 공격 가족의 관측·방어를 위한 후보이며, 모든 공격을 판별하는 정의나 보장이 아니다.
FUBA와 BadFU는 FL에서 백도어를 주입하는 효과와 이를 억제하는 효과를 함께 학습시키고,
FU에서 억제 쪽을 선택적으로 제거하여 남은 백도어를 활성화한다.
본 연구는 이를 주입과 억제의 역할이 FL에서는 함께 반영되지만 FU에서는 서로 다른 삭제 비율을
갖는 구조로 정식화한다. 공통 단위는 클라이언트 신원이 아니라 삭제 가능한 학습 기여다.
본 개정안은 원시 업데이트 이력과 실제 요청을 삭제·보존 시계열로 나누고,
시간 평활과 양의 준정부호 투영으로 상쇄 연산자를 구성한다.
이 연산자의 노출량, 정상 클래스 손실 및 FU 최소 진행량을 명시적 제약으로 두어,
실제 FU 제안에 가장 가까운 갱신을 구한다. 고정 top-k나 악성 신원 문턱은 필요하지 않다.
앞서 개발한 부호별 기여 장부는 양·음 기여의 상대 삭제율을 비교하는 기준선으로 보존한다.
국소 가법 모형에서 선택적 삭제가 숨은 효과를 노출하는 조건과
집계가 내부 장력을 가릴 수 있는 관측 한계를 분석한다.
기존 FUBA 학습 기록6개를 언러닝 모델 생성 없이 재생한 결과,
20라운드와 최종 시점에서 실제 억제자와 표적을 모두 점수1위로 찾았으나
초기 정상 후보와 표적 혼동이 남았다.
이는 구조 점수의 예비 근거이며, 정상 오탐률이나 원래 BadFU의 방어 성능 검증은 아니다.
기존 언러닝 시뮬레이션 실험은 효과 확인과 비교를 위한 부록으로 보존한다.
현재의 수식 수정만으로 방어의 신규성을 확정하지 않는다.
기존 기능 좌표 제어와 개정 연산자 제어, 실제 손실 검사 및 전체 기록 replay 실행기를 구현했다.
한 개발 모델에서 정답 트리거로 관측한 장력 상위8좌표는 모두 실제 표적을 가리켰지만,
밝기 변환에서는 이를 관측하지 못했다. 이 결과는 관측 가정을 분리해 검증할 필요를 보인다.
개정안은 균일 삭제 소거와 조건부 유일해를 가지며 수치 최적성도 검사한다.
같은 모델의 두 요청에서 제약은 충족했으나 ASR은 이력 없는 보존 제어와 같았다.
따라서 전체 FU의 방어 성능, 원래 BadFU, 독립 망각 검증과 이력의 추가 가치는 아직 미확정이다.
\end{abstract}

\section*{2026년 10월 15일: 새 업로드와 주장 범위}
이번에 확인한 자료는 CIFAR-10/ResNet-18 공격 궤적 3개와 정상 궤적 4개의
생성 및 기여 빼기 진단 결과다. 새 독립 탐지 성능표는 아직 확인하지 못했다.
K=20, 100라운드, 참여율 0.4의 공격 궤적에서 최종 ASR은 시드 522/523/524에 대해
97.13/89.30/96.53\%다. 따라서 이 종료점들을 삭제 전 잠복의 확인 사례로 쓰지 않는다.
억제자 기여를 빼면 세 모델 모두 ACC 9.93\%, ASR 0\%가 되지만,
이는 효용 붕괴를 동반하는 1차 근사 진단이며 실제 FU 방어 성능이 아니다.
공격 없는 CIFAR IID에서도 같은 종류의 정상 클라이언트 기여 빼기 후
ACC가 82.43\%에서 49.91\%로 낮아진다.
단순 변화 크기와 공격 원인을 분리해야 하는 근거다.

해당 CIFAR 검증 기준은 첫 결과를 본 뒤 완화되었으므로 현재 자동 장부의
통과 표시를 사전 고정된 잠복 공격 검증으로 해석하지 않는다.
기존 여섯 개발 궤적의 높은 순위 역시 독립 FPR/TPR 검증과 분리한다.
현재 신규성은 확정된 성과가 아니라 요청 조건부 이력의 추가 가치에 대한 가설이다.
상세 출처는 저장소의 \texttt{docs/EVIDENCE\_AUDIT.md},
원시 기록은 \texttt{FUBA/logs/trajectories/registry.json}에 보존한다.

\section{문제: FL의 잠복을 FU의 삭제로 깨뜨리는 공격}
\paragraph{변하지 않는 위협 범위.}
공격자는 FL에서 학습한 효과를 FU의 삭제 과정으로 활성화한다.
실험에서는 삭제 전·후의 ASR뿐 아니라 양쪽의 정상 과제 성능을 함께 확인한다.
장력이 큰 사례만을 공격 정의에 포함시켜 장력 검출이 성공하도록 문제를 순환 정의하지 않는다.
의도적으로 심은 백도어와 공격 없는 non-IID 모델의 유사한 오분류는 원인을 구별하여 평가하되,
실제 방어의 목표는 사용자의 의도를 알아내는 것이 아니라 망각을 이행하면서 표적 오작동을 억제하는 것이다.

일반적인 백도어 방어는 학습 또는 추론 중 드러난 공격 효과를 검사한다.
여기서 다루는 공격은 두 단계의 연결을 악용한다.
FL 동안에는 주입 효과와 억제 효과를 함께 반영하여 행동을 숨기고,
FU에서는 억제에 해당하는 기여를 제거하면서 주입 효과를 남긴다.
따라서 학습 종료 모델의 겉보기 정상성만으로 이후 삭제의 안전성을 판단하기 어렵다.

FUBA의 Adv-attacker와 Adv-defender는 이 역할을 서로 다른 클라이언트에서 수행한다
\cite{fuba}. BadFU의 backdoor 표본과 camouflage 표본은 같은 악성 클라이언트 안에서
주입과 억제 역할을 수행하며, camouflage 표본 삭제로 백도어가 활성화된다\cite{badfu}.
이는 공통 구조의 서로 다른 실현이다. 단위 해상도의 차이를 공격 구조의 차이로 취급하지 않는다.

이 위험은 언러닝 근사 오차가 있어야만 성립하는 것이 아니다.
억제 표본을 정확히 삭제해도 보존 데이터에 남은 백도어 효과가 학습될 수 있으므로,
망각의 정확성과 백도어 안전성은 별도의 조건이다.
공격 논문들이 설명한 이 기제를 하나의 이름으로 묶는 것만을 신규성으로 삼지는 않는다.
이 구조에서 삭제 이전에 계산할 수 있는 방어 절차를 도출하고 검증하는 것이 목표다.

\begin{table}[ht]\centering\small
\caption{같은 FL$\to$FU 구조에서의 역할 대응. 클라이언트 수는 구조의 정의가 아니다.}
\begin{tabularx}{\linewidth}{lXXX}\toprule
공격 & FL의 주입 역할 & FL의 억제 역할 & FU에서 제거하는 부분\\\midrule
FUBA & Adv-attacker 기여 & Adv-defender 기여 & 억제자의 학습 기여\\
BadFU & backdoor 표본 기여 & camouflage 표본 기여 & camouflage 표본 기여\\\bottomrule
\end{tabularx}
\end{table}

본 연구의 중심 질문은 다음과 같다.
\begin{quote}
FL에서 서로 가려진 기능적 기여를 장력으로 기록하고,
FU 요청이 그중 억제 쪽을 더 제거하는 방향을 찾아 실제 보정 갱신에서 제어할 수 있는가?
\end{quote}
우리는 \emph{삭제 경계로 조건화한 장력 장부}를 연구 후보로 제시한다.
단순한 기울기 충돌, 반대 정렬, 또는 삭제 후 ASR 상승 자체를 새 발견으로 주장하지 않는다.
기여의 대상은 학습 중의 관계를 삭제 가능한 단위에 연결하고,
요청이 그 관계를 보존하는지 깨뜨리는지를 검사하는 방어 절차다.
정상 충돌과 악성 은닉을 구별하는 검출력은 별도 실험으로 입증해야 한다.
``FL 이력을 FU에 쓴다''는 것 역시 단독 신규성이 아니다.
UnlearnGuard는 이미 이력 기반 FU 갱신 예측과 검사를 사용한다\cite{unlearnguard}.
우리가 검증해야 하는 차이는 FL에서 함께 있을 때 가려졌던 기여가 삭제로 분리되는 관계가
개별 갱신의 크기·방향·시간적 일관성보다 추가 정보를 주는지에 있다.

\section{관련 연구와 논리 전개}
MASA는 개입으로 행동 차이를 관측하고 non-IID 혼동을 완화하는 절차를 둔다\cite{masa}.
FLDetector는 이력에서 비교 기준을 도출해 예상 갱신과의 차이를 검사한다\cite{fldetector}.
여기서 가져오는 것은 \emph{관측 가능한 현상에서 측정과 판정을 도출하는 방식}이다.
본 연구에서 그 현상은 삭제 후 변화가 아니라 삭제 이전에 유지되는 주입--억제 관계다.

정상 다중 과제에서도 기울기 충돌이 생길 수 있다\cite{pcgrad}.
따라서 ``장력이 있으면 공격''이라는 충분조건은 사용하지 않는다.
FLAME과 DeepSight가 탐지와 잔여 영향 처리를 구분하듯\cite{flame,deepsight},
장력 경보, 공격자 귀속, 실제 복구, 망각 완료도 각각 검증해야 한다.
언러닝 과정의 갱신을 다루는 UnlearnGuard\cite{unlearnguard}와는 실행 시점 및 관측 정보를 맞춰 비교한다.
선행 연구의 존재만으로 본 알고리즘의 신규성이 확정되거나 부정되는 것은 아니다.

\subsection{다른 방어들은 무엇을 신규성으로 연결했는가}
다음은 원논문의 방법과 기여 진술을 읽고 도출한 설계 비교다.
어떤 논문이든 non-IID에서 공격과 정상 의도를 항상 구별했다는 의미는 아니다.
\begin{table}[ht]\centering\small
\caption{관측의 공백에서 알고리즘 기여로 이어지는 전개. 우리 설계로의 적용은 해석이다.}
\begin{tabularx}{\linewidth}{lXXX}\toprule
연구 & 기존 방식의 공백 & 새로운 관측 또는 조치 & 우리에게 필요한 검증\\\midrule
MASA\cite{masa} & 정적 갱신 통계에 가려진 백도어 & 개별 언러닝 동역학을 측정하고 사전 모델 융합으로 non-IID 혼동을 완화 & 새 관측이 정상 이질성에도 유지되는가\\
BackdoorIndicator\cite{indicator} & 작은 악성 갱신과 정상 갱신의 중첩 & OOD 지표 과제를 심어 능동 측정하고 BN 통계 이동을 통제 & 기존 관측에 없던 정보가 실제 생기는가\\
FedREDefense\cite{fedredefense} & 클라이언트 사이 분포 비교의 불안정성 & 각 갱신의 학습 재구성 오차를 사용 & 실제 poisoned-data 학습도 재구성되면 남는 신호가 있는가\\
DeepSight\cite{deepsight} & 단순 이상치가 정상 non-IID까지 배제 & 모델 내부의 여러 표현과 군집 검사, 잔여 영향 clipping & 희귀 클래스 기여를 보존하며 위험 영향만 줄이는가\\
FLAME\cite{flame} & 필터 누락과 과도한 잡음의 정확도 손실 & 군집·clipping으로 필요한 잡음량을 줄이는 결합 & 각 구성요소가 다음 단계의 비용이나 실패를 줄이는가\\
UnlearnGuard\cite{unlearnguard} & FU 도중 새로 제출되는 악성 갱신 & FL 이력으로 FU 갱신을 예측하고 크기·방향을 검사 & 정상처럼 보였던 FL 기록 자체의 잠복 관계에도 충분한가\\\bottomrule
\end{tabularx}
\end{table}
특히 FedREDefense의 재구성 가능성을 악의 판별과 동일시하지 않는다.
BadFU처럼 오염 데이터를 사용해 실제로 최적화한 갱신이 재구성될 가능성은 별도 검증해야 한다.
BackdoorIndicator의 OOD 및 지표 과제 가정도 원래 BadFU/FUBA에 자동으로 성립하지 않는다.
논리 전개에서 가져올 것은 새 관측의 근거, 혼동 통제, 조치와 소거의 연결이다.

\subsection{FedEraser의 기록을 사용할 때의 정확한 전제}
FedEraser는 서버가 보관한 클라이언트별 과거 모델 업데이트를 재사용한다\cite{federaser}.
각 minibatch의 원시 gradient 전체를 별도로 요구하는 것은 아니다.
과거 업데이트 $u_j^t$와 남은 데이터로 짧게 보정 학습한 업데이트 $\widehat u_j^t$로
\begin{equation}\label{eq:federaser-original}
 \widetilde u_j^t=\|u_j^t\|\frac{\widehat u_j^t}{\|\widehat u_j^t\|}
\end{equation}
를 구성한 뒤 보정 갱신을 집계한다. 영벡터 처리와 참여·가중 규약은 구현에서 지정해야 한다.
따라서 같은 저장 조건의 FUBA 클라이언트 수준 감사에서는 기존 이력을 활용할 수 있다.
서버의 원시 데이터 접근이나 표본별 기여 장부를 공짜로 얻는다는 뜻은 아니다.
원래 FedEraser는 클라이언트 제거 방법이므로 BadFU의 표본 삭제를 지원하는 변형의 정보와
보정 데이터 조건은 별도로 명시한다.

로컬 \texttt{federaser\_faithful.py}는 현재 클라이언트 평균을 먼저 만든 뒤 층별 노름을
적용한다. 일반적으로 식~\eqref{eq:federaser-original}의 클라이언트별 보정 후 집계와 다르므로,
이 파일 이름만으로 원논문 정합성을 인정하지 않는다. 후속 FU 비교 전에 이 연산 순서와
부분 참여 처리를 검증해야 한다. 현재 장력 캐시 재생은 이 FU 구현을 호출하지 않는다.

\section{공통 모형: 역할과 삭제 단위를 분리한다}
\subsection{공격 가족의 정의}
전체 학습 단위를 $\mathcal U$라 하고, 그 안의 주입 집합 $A$와 억제 집합 $S$를 생각한다.
이 둘은 같은 클라이언트에 속할 수 있다. 입력 영역 $X^*$와 표적 행동에 대해
공격자는 FL 종료 시 낮은 공격 효과와 FU 후 높은 공격 효과를 함께 목표로 한다.
이 주입--억제 모형은 FL$\to$FU 잠복 공격을 분석하는 중심 기제 가설이다.
잠복--활성화의 행동 정의만으로 모든 비선형 공격에 가법 상쇄가 필수라는 정리를 주장하지 않는다.

누적 요청은 $r\in[0,1]^{|\mathcal U|}$이며 $r_u$는 단위 $u$의 삭제 비율이다.
이는 기록된 기능 기여를 제거하는 비율이다. 표본 수의20\% 삭제를 기여의20\% 제거로
자동 치환하지 않는다. 임의 표본 삭제에는 그 경계를 표현하는 더 세밀한 단위가 필요하며,
원자적 단위의 실제 삭제는 $r_u\in\{0,1\}$로 표현한다.
삭제 경계란 $r$이 정하는 제거 부분과 보존 부분의 구분이다.
FL 동안 함께 반영된 역할들이 이 경계에서 서로 다른 비율을 갖는 것이 핵심이다.
실제 언러닝 연산자 $\mathcal R$와 그 뒤의 모델은 위협의 결과를 정의하지만,
아래 장력 점수를 계산할 때는 $\mathcal R$을 호출하지 않는다.

\subsection{기능 기여를 기록하는 방법}
같은 입력 $x$와 출력 방향 $c$에서 기능 좌표 $m_c(x;\theta)$를 정한다.
예를 들어 로짓 $f_c-\frac1C\sum_df_d$를 사용할 수 있다.
라운드 $t$의 실제 가중 학습 갱신 $u_u^t$에 대해
\begin{equation}\label{eq:train-effect}
 e_u^t(x,c)=\nabla_\theta m_c(x;\theta_{t-1})^\top u_u^t,\qquad
 E_u^t(x,c)=\sum_{s\in W_t}e_u^s(x,c)
\end{equation}
를 정의한다. $W_t$는 정해진 기록 창이며 누적 장부이면 $W_t=\{1,\ldots,t\}$이다.
JVP는 현재 학습 모델의 방향 미분이며 삭제 모델을 만들지 않는다.
기존 CCD 구현은 각 라운드의 실제 가중 갱신을 더한 모델의 로그확률 차이를 측정한다.
이 역시 과거 기여를 제거하는 언러닝은 아니지만 추가 순전파는 필요하다.
기존 저장 캐시는 그 효과를 입력 평균으로 줄인 클래스 벡터다.

식~\eqref{eq:train-effect}의 누적량은 기능 기여의 근사 장부다.
시간에 따른 모델 변화, 최적화 상태, 비선형 상호작용 때문에 정확한 삭제 영향과 동일하지 않다.
입력별 상쇄가 평균에서 사라질 수 있으므로 새 설계는 입력 좌표를 먼저 보존한다.
공개 입력 또는 정당하게 접근 가능한 요청 입력을 사용하며 공격 표적과 트리거 정답을 제공받는다고 가정하지 않는다.

\subsection{BadFU를 포함하는 데 필요한 기록 해상도}
BadFU에서도 동일한 $A,S,r$ 정의를 사용한다.
다만 같은 클라이언트 내부의 두 역할을 관측하려면 데이터 묶음별 기록 등
요청을 표현할 수 있는 해상도의 정보가 필요하다.
서버가 클라이언트의 합산 갱신만 받는다면 그 합을 임의의 표본 기여로 정확히 분해할 수 없다.
이는 두 공격의 구조가 다르다는 뜻이 아니라, 방어가 보존해야 할 정보의 해상도를 정하는 조건이다.
신뢰 가능한 묶음 기록과 검증 절차를 갖추는 방법은 현재 구현에서 추가로 필요한 부분이다.

\section{기능적 장력에서 삭제 경계 점수로}
이 절에서는 $(t,x,c)$를 고정하고 $E_u=E_u^t(x,c)$로 줄여 쓴다.
양의 밀기와 음의 밀기를
\begin{equation}
 P=\sum_u[E_u]_+,\qquad N=\sum_u[-E_u]_+
\end{equation}
로 분리한다. 순효과는 $P-N$, 총 기여량은 $P+N$이다.
장력의 크기와 상쇄 비율은
\begin{equation}\label{eq:tension}
 T=P+N-|P-N|=2\min(P,N),\qquad
 H=\begin{cases}T/(P+N),&P+N>0,\\0,&P+N=0.\end{cases}
\end{equation}
이다. $H$만 크면 미세한 잡음도 완전 상쇄처럼 보일 수 있으므로 크기를 함께 유지한다.
여기까지는 FL 내부의 장력이며 어떤 삭제가 위험한지는 구별하지 못한다.

요청 $r$을 기존 장부에 적용하여 네 양을 계산한다.
\begin{align}
 P_D&=\sum_u r_u[E_u]_+,&N_D&=\sum_u r_u[-E_u]_+,\\
 P_K&=\sum_u(1-r_u)[E_u]_+,&N_K&=\sum_u(1-r_u)[-E_u]_+.
\end{align}
$D$는 제거 부분, $K$는 보존 부분이다. 이전 후보(v2)는
\begin{equation}
 B_r^{\mathrm{v2}}=H[\min(N_D,P_K)-\min(P_D,N_K)]_+
\end{equation}
였다. 그러나 $(E_A,E_S)=(10,-8)$을 각각80\% 제거하면 보존 순효과가2에서0.4로
줄어들고 상대 균형은 같은데도 $B_r^{\mathrm{v2}}=16/45>0$이다.
양방향 질량의 차만으로는 불완전 상쇄에서 동일 비율 삭제를 배제하지 못한다.
따라서 현재 후보(v3)는 효과 질량으로 가중한 삭제율을 사용한다.
\begin{equation}
 d_-(r)=N_D/N,\qquad d_+(r)=P_D/P,\qquad M=\min(P,N).
\end{equation}
$P=0$ 또는 $N=0$이면 점수는0으로 정의한다. 그 외에는
\begin{equation}\label{eq:boundary}
 D_r=[d_-(r)-d_+(r)]_+,\qquad B_r=H M D_r
\end{equation}
로 정의한다. $M$은 서로 가려질 수 있는 질량, $H$는 사전 상쇄 비율,
$D_r$는 억제 쪽을 더 제거하는 정도다. 클라이언트 수나 표본 수의 비율이 아니다.
이는 실제 인과적 짝짓기를 복원한 결과가 아니라 부호 장부의 선택적 삭제 통계다.

\begin{table}[ht]\centering\small
\caption{$E_A=+10,E_S=-10$일 때 장력 $T=20$은 같지만 삭제 경계 점수는 다르다.}
\begin{tabular}{lrrr}\toprule
요청 & $(r_A,r_S)$ & $B_r$ & 장부의 보존 순효과\\\midrule
아무것도 제거하지 않음 & $(0,0)$ & 0 & 0\\
억제만 제거 & $(0,1)$ & 10 & 10\\
주입만 제거 & $(1,0)$ & 0 & $-10$\\
둘 다 제거 & $(1,1)$ & 0 & 0\\
둘 다20\% 제거 & $(.2,.2)$ & 0 & 0\\
억제만20\% 제거 & $(0,.2)$ & 2 & 2\\\bottomrule
\end{tabular}
\end{table}
표의 수치는 실제 FU 결과가 아니라 구조 점수의 산술 예다.
장력의 존재와 그 장력을 깨는 요청을 구분하는 것이 이 설계의 핵심이다.

\section{분석: 선택적 제거가 왜 중요한가}
\begin{proposition}[두 역할의 삭제 비대칭]
한 기능 좌표에서 주입 효과 $b\ge\beta>0$, 억제 효과 $s=-b+\eta$, $|\eta|\le\varepsilon$라 하자.
두 역할의 삭제 비율이 $r_A,r_S$이면 가법 장부의 보존 효과는
\begin{equation}\label{eq:asymmetry}
 (1-r_A)b+(1-r_S)s=(r_S-r_A)b+(1-r_S)\eta.
\end{equation}
특히 $r_S\ge r_A$일 때, 실제 FU 효과와 이 근사 사이 오차의 절댓값이 $\delta$ 이하이면
실제 기능 효과의 하한은 $(r_S-r_A)\beta-(1-r_S)\varepsilon-\delta$다.
\end{proposition}
\begin{proof}
$s=-b+\eta$를 대입하고 $b\ge\beta$, $|\eta|\le\varepsilon$와 오차 경계를 적용한다.
\end{proof}
이는 가법 근사와 오차 한계 아래의 조건부 명제다.
관측하지 못한 트리거 입력의 ASR이나 임의 언러닝 연산자의 성능을 보장하지 않는다.
동일 비율 삭제는 큰 주입 항을 소거하지만 억제를 더 제거하면 그 차이에 비례한 항이 남는다.
두 역할이 같은 클라이언트인지 여부는 식에 등장하지 않는다.

\begin{proposition}[비율 보존과 완전 균형에서의 단순화]
$P,N>0$이면 $0\le B_r\le H\min(P,N)$이다.
모든 단위의 삭제 비율이 같으면 $P=N$ 여부와 무관하게 $B_r=0$이다.
한쪽 부호 기여를 같은 삭제 비율을 가진 여러 단위로 나누어도 점수는 같다.
특히 $P=N>0$이면
\begin{equation}
 B_r=[N_D-P_D]_+=[P_K-N_K]_+.
\end{equation}
\end{proposition}
\begin{proof}
$d_-,d_+\in[0,1]$이고 동일 비율 제거에서는 둘 다 그 비율과 같다.
동일 부호·동일 삭제율의 분할은 $P,N,P_D,N_D$를 보존한다.
$P=N$이면 $H=1$, $M=P$를 식~\eqref{eq:boundary}에 대입하면 된다.
\end{proof}
이 특수 경우 점수는 장부의 보존 방향 효과로 단순화된다.
따라서 식 하나를 새로운 이론으로 내세우는 것은 충분하지 않다.
입력·시간·삭제 단위를 보존하는 기록과 정상 대조에서의 검증이 알고리즘 기여에 필요하다.

\begin{proposition}[집계 해상도의 한계]
관측이 같은 단위의 합 $e_1+e_2$뿐이라면,
$(e_1,e_2)=(0,0)$과 $(b,-b)$는 같은 관측0을 만든다.
그 합만 사용하는 검사는 내부 장력의 존재나 두 번째 부분만 제거했을 때의 효과를 항상 구별할 수 없다.
\end{proposition}
\begin{proof}
두 경우의 관측은 같지만 장력은 각각0과 $2|b|$이고, 두 번째 부분 제거 후 남는 효과는 각각0과 $b$다.
\end{proof}
이 명제는 BadFU를 제외하는 근거가 아니라, 클라이언트 내부 삭제까지 방어하려면
그 경계를 표현할 수 있는 기록이 필요하다는 설계 근거다.
또한 정상 목적의 충돌도 동일한 장력 수치를 만들 수 있으므로 장력은 악의의 충분조건이 아니다.

\section{알고리즘과 구현 범위}
\begin{enumerate}
\item \textbf{FL 기록.} 가중 학습 기여의 기능 효과를 동일 입력·출력 좌표에서 기록하고,
      클라이언트 또는 데이터 묶음의 삭제 식별자에 연결한다. 평균 전에 입력 좌표를 보존한다.
\item \textbf{잠복 관계 감시.} 정해진 시간 창에서 $T,H$와 기여의 부호를 누적한다.
      같은 영역과 출력 방향에서 지속되는 큰 장력을 후보로 남긴다. 단일 주입자를 배제하지 않는다.
\item \textbf{FU 경계 검사.} 요청의 누적 $r$로 $B_r$를 계산한다.
      요청 전에는 가능한 단위 경계의 후보 점수를 기록할 수 있다.
      어느 경우에도 후보 언러닝 모델을 생성하지 않는다.
\item \textbf{정상 대조와 판정.} 데이터 이질성·참여량·삭제량을 고려한 정상 학습 에피소드에서
      입력·클래스·시간·요청 선택까지 포함한 전체 판정을 보정한다.
      아직 이 보정을 완료하지 않았으므로 새 점수는 순위와 진단 증거로 보고한다.
\item \textbf{대응과 검증.} 위험 후보의 근거를 검토한 뒤 전체 요청을 이행하는 복구를 구성하고,
      보존·망각·잔여 공격 및 후속 재활성화를 검증한다. 억제 성분을 되넣는 것은 해결로 세지 않는다.
\end{enumerate}

현재 \texttt{tension\_audit.py}는 식~\eqref{eq:tension}--\eqref{eq:boundary}를 구현하며
삭제 가능한 단위 이름이 클라이언트인지 같은 클라이언트의 묶음인지에 의존하지 않는다.
\texttt{replay\_tension.py}는 저장된 학습 장부만 읽고 점수를 계산한다.
\texttt{collect\_tension.py}는 기존 CCD의 양의 학습 갱신 관측을 재사용하면서
입력별 로그확률 및 중심화 로짓 효과를 보존한다. 이는 JVP가 아닌 유한 갱신 관측이다.
\texttt{experiment\_tension.py}는 정상/공격 학습, 장부 수집, 소거 및 정상 에피소드 보정을 연결한다.
입력별 JVP, 신뢰 가능한 묶음 기록 및 새 점수로 실제 복구하는 통합은 남은 구현이다.
기존 CCD는 학습 중 장력 관측의 출발점으로 유지하되, 양의 다리2개 이상 같은 FUBA 설정 의존 규칙을
새 공통 점수의 정의에 넣지 않는다.
기존 DOSE·soft 반응 검사와 공동삭제는 비교군 및 사후 효과 확인에 사용한다.

\section{실측: 학습 기록만으로 무엇이 남는가}
\subsection{새 구조 점수의 캐시 재생}
R-034는 기존 FUBA IID/owner 각3개, 총6개 학습 캐시를 사용한 v3 재생이다.
이전 v2의 R-032--R-033은 보존하며 서로 다른 버전의 수치를 섞지 않는다.
새 모델 학습, 언러닝 모델 생성, 실제 FU 호출은0회다.
학습 시점10/20/40/99에서 각 클라이언트 단독 삭제 경계의 점수를 계산하고
클래스 최대값으로 요청 후보를 정렬했다. 역할과 표적 정답은 정렬 후 채점에만 쓴다.
이는 실제 요청을 실행한 것이 아니라 이미 측정된 장부와 요청 메타데이터의 연산이다.

\begin{table}[ht]\centering\small
\caption{새 경계 점수의 기술적 순위 분석. 문턱을 보정한 탐지율이 아니다.}
\begin{tabular}{lrr}\toprule
학습 시점 & 실제 억제자4가 점수1위 & 1위 요청자와 표적8 모두 일치\\\midrule
10 & 5/6 & 4/6\\
20 & 6/6 & 6/6\\
40 & 6/6 & 6/6\\
99 & 6/6 & 6/6\\\bottomrule
\end{tabular}
\end{table}
IID s0의10시점은 정상12가1위였다.
owner s0의10시점은 요청자4를 찾았지만 클래스3을 택했다.
따라서 큰 장력의 요청 순위와 실제 공격 표적 식별은 별도 문제다.
이 결과를 정상 오탐률이나 새로운 조기 경보 라운드로 보고하지 않는다.
캐시는 입력 평균 로그확률 효과이므로 입력별 JVP 명세를 전부 검증한 실험도 아니다.
같은 단독 삭제 비교에서 단순 기준 $L_r=[N_D-P_D]_+$의 요청자·표적 동시 일치는
20시점0/6, 최종4/6이었다. v2는 두 시점 모두5/6이었다.
이는 이미 관찰한 개발 궤적의 기술적 차이이며 미관측 공격이나 정상 대조에서의 우월성을 입증하지 않는다.
장력 단독은 요청과 무관한 값이므로 요청 순위의 동점 해소를 검출력으로 해석하지 않는다.

\subsection{기존 근거의 보존과 정확한 위치}
기존 CCD의 R-013 및 캐시 재생 기록은6개 FUBA 궤적에서11--16라운드 경보를 보인다.
이는 기존 FUBA 설정의 장력 관측 근거로 남긴다.
다리 수와 최대 점유율 규칙이 포함되어 있으므로 이를 BadFU까지의 공통 검출 성능으로 전환하지 않는다.
BadFU 변형의 공개 클래스 평균 장부에서 경보가 없었다는 기존 음성 결과는
해당 관측기의 실패이며, BadFU에 주입--억제 구조가 없다는 증거가 아니다.

이번 장력 단위 검사15개는 같은 클라이언트 내부 역할과 클라이언트 간 역할의 동일성,
동일 비율 제거, 여러 주입자/억제자, 입력 평균의 정보 손실, 정상 반대 기여 반례를 포함한다.
보정과 데이터 분리 검사7개 및 기존 검사를 포함한 전체69개가 통과했다.
R-035의 실제 저장 모델 스모크는2개 전이와 입력16개에서 수집기를 확인했으며,
FedAvg 갱신 복원 상대 오차의 최댓값은 $3.95\times10^{-7}$ 미만이었다.
이 값은 언러닝 영향 근사 오차가 아니다. 동작 검사는 현실 방어 성능 실험과 구분한다.
이전84조건의 언러닝 후 soft 감사와 공동삭제 실측은 부록에 보존한다.
그 결과를 본 장력 알고리즘의 성능으로 합산하지 않는다.

\section{실제로 사용할 때의 실패 지점과 수정 원칙}
\subsection{정상 충돌과 선택 편향}
정상 non-IID 학습도 큰 $T,H$와 선택적 삭제를 만들 수 있다.
시간·출력·요청을 많이 검사한 뒤 가장 큰 점수만 고르면 고정 문턱의 오탐이 증가한다.
독립적인 정상 학습 에피소드 $i$에 대해 사전에 고정한 요청·시점 집합의 최대 점수
\begin{equation}
 Z_i=\max_{t,r,c}\frac1{|X|}\sum_{x\in X}B^{(i),t}_r(x,c)
\end{equation}
를 만들고 동일한 분포 조건의 정상 에피소드끼리 보정한다.
입력 평균을 먼저 취하는 소거군은 별도 설정으로 취급한다.
정상 보정 수가 $n$이면 $k=\lceil(n+1)(1-\alpha)\rceil$로 두고,
$k\le n$일 때 $\tau=Z_{(k)}$, 아니면 $\tau=+\infty$를 사용한다.
검사는 $Z>\tau$일 때만 경보를 낸다. 교환 가능성과 미리 고정된 점수 가족 아래에서
이는 새 정상 에피소드의 경보 확률을 $\alpha$ 이하로 제한하는 순위 보정이다.
정상19개는 $\alpha=.05$에서 유한 문턱이 가능한 최소치일 뿐 충분한 검출력이나
정확한 경험적 오탐률 추정의 보장이 아니다. 여러 소거 중 시험 성능이 가장 좋은 것을
선택하면 개별 보정을 최종 방법의 보장으로 전환할 수 없다.

\subsection{기록을 믿을 수 있는가, 삭제 경계를 표현할 수 있는가}
서버가 검증할 수 없는 악성 클라이언트의 자체 묶음 장부는 위조될 수 있다.
묶음의 사전 식별자 고정은 사후 재분류를 막지만 기록의 진실성을 증명하지 않는다.
BadFU 실험에는 원래 같은 클라이언트 안의 두 역할을 유지하고,
요청과 연결 가능한 기록을 독립적으로 계측하는 조건이 필요하다.
역할 정답으로 만든 oracle 묶음은 관측 가능성 상한이며 실제 서버 방어로 보고하지 않는다.
신뢰 가능한 기록이 없거나 요청이 묶음 내부를 가르면 ``점수0이므로 안전'' 대신
관측 불충분으로 반환해야 한다. 보안 집계로 개별 기여가 보이지 않는 설정도 별도 프로토콜이 필요하다.

\subsection{함수 관측의 누락과 시간 변화}
공개 입력에서 효과가 전혀 나타나지 않으면 입력별 기록도 이를 복원하지 못한다.
입력 평균 전·후 점수, 로그확률·중심화 로짓, 누적·최근10/20전이를 비교한다.
서로 다른 모델 시점의 국소 효과를 누적하는 오차는 창 크기 실험으로 진단할 수 있으나
창을 짧게 만드는 것만으로 해결되지는 않는다. 짧은 창은 초기 주입도 잊을 수 있다.
같은 클라이언트의 시간별 반대 효과도 단위 내 합산으로 사라질 수 있다.
이를 발견했다고 삭제할 수 없는 시간 조각을 독립 삭제 단위로 가정해서는 안 된다.

\subsection{구조 경보에서 복구로 넘어가는 간극}
높은 $B_r$는 위험한 삭제 관계의 후보이지 악성 사용자 귀속이나 망각 완료 증명은 아니다.
재억제로 ASR만 낮추는 대응은 제거한 역할을 복원할 수 있으므로 목표가 아니다.
후속 실험은 전체 요청을 반영한 재학습 또는 기존 FU에 후보 처리 절차를 붙이고,
보존 정확도·망각·ASR·후속 학습의 재활성화를 함께 검사해야 한다.
새 점수가 해당 복구를 수행했다는 주장은 현재 보류한다.

\section{신규성의 검증과 필요한 실험}
논문의 기여 주장은 \emph{FL의 잠복 관계를 FU의 삭제 경계에 연결하는 방어}에 둔다.
다음 검증이 이 주장을 지지하거나 반박해야 한다.
먼저 공개한 실행 묶음은 FUBA IID/owner 각 정상 보정3개·정상 시험3개·기존 공격3개인
18개 예비 에피소드다. 정상 보정이 부족하므로5\% 경보 성능은 보고하지 않는다.
확인 묶음은 조건별 정상 보정19개·정상 시험10개·새 공격5개, 총68개다.
주 비교는 누적 로그확률 클래스 평균의 v3와 $L_r=[N_D-P_D]_+$이며,
v2, 장력 단독, 제거된 음의 기여량 및 입력·시간·좌표 소거를 함께 저장한다.
이 실행 묶음에는 원래 BadFU 표본 기록과 실제 FU·복구가 아직 포함되지 않는다.
\begin{enumerate}
\item \textbf{경계의 필요성:} 같은 FL 장력에서 억제만/주입만/둘 다/동일 비율/분할 삭제를 비교한다.
      지금은 수식과 단위 검사이며, 실제 FU로 예측과 결과를 비교해야 한다.
\item \textbf{두 공격의 공통성:} FUBA의 클라이언트 역할과 원래 BadFU의 표본 역할에 같은 점수를 적용한다.
      BadFU를 별도 공격 부류로 바꾸거나 클라이언트 분리 변형만으로 검증을 대신하지 않는다.
\item \textbf{기록 해상도:} 같은 모델에서 클라이언트 합산/묶음 기록, 클래스 평균/입력별 기록을 소거한다.
      정상 대조를 포함해 어느 정보가 실제로 필요한지 측정한다.
\item \textbf{정상 충돌 통제:} 공격 없는 동일 K·학습 길이의 IID/Dirichlet/owner 조건에서
      학습 시드를 나누어 보정하고, 장력 단독과 경계 조건부 점수를 동일 FPR에서 비교한다.
\item \textbf{연산자와 복구:} 근사 감산/FedEraser/재학습으로 장부 예측의 오차를 측정하고,
      추정 후보로 복구한 뒤 ACC·망각·ASR·후속 재활성화를 함께 평가한다.
\item \textbf{적응형 공격:} 측정 규칙을 아는 공격자가 주입과 억제를 함께 재학습하는 조건에서
      잠복성, 활성화, 장력 점수의 관계를 확인한다. 저장 업데이트의 배율 변경만으로 대체하지 않는다.
\end{enumerate}
언러닝 시뮬레이션은 이 실험에서 장부가 예측한 효과를 확인하는 도구이며,
검출 알고리즘의 핵심 단계 또는 신규성 자체로 제시하지 않는다.
동일 정상 경보 수준에서 단순 기준보다 이득이 없다면 복잡한 점수를 핵심 기여에서 내린다.
BadFU의 검증 가능한 묶음 기록에서 신호를 확보하지 못하면 두 공격의 구조적 공통성은 유지하되
두 공격을 실제 방어했다는 실험 주장은 철회한다.

\section{실제 FU 갱신 제어: 구현 명세와 검증할 주장}
이 절의 방향 선택·국소 제어·클라이언트별 보정 코어를 구현했다.
실행기는 기존 MNIST 궤적에서 실제 보정 한 단계의 소거 실험을 수행한다.
전체 FU 파이프라인, 원래 BadFU의 표본 삭제 어댑터 또는 보안·망각 인증이 완성된 것은 아니다.

\subsection{먼저 분리해야 할 세 가지 명제}
첫째, 주입 $b\ge\beta>0$와 억제 $s=-b+\eta$, $|\eta|\le\varepsilon<\beta$가
실제 측정 좌표에 존재하면 $T=2\min(b,-s)\ge2(\beta-\varepsilon)$다.
이 의미에서 강한 은닉 기여는 장력을 요구한다.
그러나 높은 ASR 증가율만으로 공개 입력에서 큰 절대 장력을 보장할 수는 없다.
여러 입력의 분류 마진이0에 가까우면 아주 작은 변화로도 ASR이 크게 바뀌며,
로짓의 양의 상수 배율 변경은 예측을 보존하면서 로짓 장력의 크기를 바꾼다.

둘째, 정상/공격의 관측 $O$가 동일하면 $O$만으로 의도를 항상 구별할 수 없다.
이는 모든 non-IID 설정에서 모든 방어가 불가능하다는 명제가 아니다.
추가 입력, 과제 지식, 과정 관측 또는 조치의 목표를 바꾸면 비교할 정보가 달라진다.
셋째, 공격 탐지가 불가능한 특정 경우에도 정당한 입력 의미와 보존 과제가 주어지면
관측된 표적 오작동을 완화하는 목표를 세울 수 있다. 그 외부 기준의 비용과 적용 범위를 숨기지 않는다.

\subsection{정상 클래스 편향을 제거한 기능 좌표 후보}
작업 의미를 유지한다고 사전에 정한 입력 변환 집합 $\mathcal A$에 대해
\begin{equation}\label{eq:paired-coordinate}
 q_c(x,a;\theta)=m_c(a(x);\theta)-m_c(x;\theta),\qquad a\in\mathcal A,
\end{equation}
여기서 $m_c=f_c-C^{-1}\sum_df_d$인 중심화 로짓을 후보 좌표로 둔다.
이 좌표에서 각 실제 학습 갱신의 효과를 기록하고,
입력 평균 전에 장력을 계산하여 같은 누적 삭제율 $r$에 조건화한다.
정상 클래스별 기준 분포와 희귀 클래스 성능을 별도로 유지한다.
일률적으로 클래스 $c$의 로짓을 올리는 항 $b_c(\theta)$는 두 입력의 차에서 정확히 소거된다.
이는 제한된 편향 제거 성질일 뿐 일반적인 non-IID 효과 제거 정리는 아니다.
정상 특징과 변환의 상호작용은 남을 수 있으며 로짓 배율도 별도 정규화·보정 문제다.
중심화는 모든 출력에 공통인 로짓 이동도 제거한다. 클래스 마진을 사용하는 소거도 가능하다.

새 관측은 실제 트리거 정답을 탐지기에 주지 않는다.
사전에 정한 변환이 공격 특징을 전혀 자극하지 못하면 신호가 없을 수 있다.
실제 트리거를 쓰는 결과는 평가용 관측 상한으로 분리한다.
이 가정 아래서만 ``원시 장력보다 정상 non-IID 혼동을 줄인다''는 가설을 실험한다.

\subsection{과거 장력과 실제 갱신의 일치}
삭제 요청을 반영한 실제 보정 업데이트를 $\widehat u_j$라 할 때 FedEraser의 제안은
\begin{equation}
 v=\sum_{j\notin R}\bar w_j\frac{\|u_j^{\mathrm{FL}}\|}{\|\widehat u_j\|}\widehat u_j.
\end{equation}
클라이언트별 전체 벡터를 먼저 보정하고 집계한다. 새 코어는 매 라운드의 참여자를 명시적으로
받으며 삭제한 단위에 보정 데이터를 요청하지 않는다. 기존 평균 후 층별 보정 파일은 과거 변형으로 보존한다.
과거 노름이 양수인데 새 보정 방향이0이면 오류를 보고하고 옛 방향으로 대체하지 않는다.

앞서 정의한 삭제 조건부 v3 점수를 $B_p(r)$로 쓴다. 이미 계산한 실제 제안에 대한
한 번의 기능 평가로 현재 증가 $A_p$를 얻고 다음 순위를 사용한다.
\begin{equation}\label{eq:action-selection}
 A_p(v)=[q_p(\theta+v)-q_p(\theta)]_+,\quad
 S_p(r,v)=B_p(r)A_p(v),\quad
 \mathcal P_r=\operatorname{TopM}_{S_p>0}S_p.
\end{equation}
이는 실제 FU 한 제안의 평가이며 가상 삭제 요청별 FU 모델을 여러 개 만드는 탐지가 아니다.
과거 장력만으로 방향을 고르는 소거군을 유지한다. 현재 구현의 양의 top-$M$은 순위 휴리스틱이며
보정된 정상 경보 문턱이 아니다. 로짓 배율의 정규화와 독립 정상 보정은 추가 작업이다.

\subsection{실제 FU 갱신의 선택적 제어}
삭제 요청을 이행하는 FU가 제안한 갱신을 $v$라 하자.
FL 장부와 현재 삭제 경계가 선택한 소수의 위험 좌표 집합을 $\mathcal P_r$로 두고,
현재 모델에서 $g_p=\nabla_\theta q_p$, 보존 과제 proxy의 클래스별 손실 기울기를
$h_y=\nabla_\theta L_y^{\mathrm{keep}}$로 계산한다. 다음을 후보 제어 문제로 둔다.
\begin{align}\label{eq:controlled-step}
 d^*=\arg\min_d\quad&\tfrac12\|d-v\|_2^2
       +\lambda\sum_{p\in\mathcal P_r}[g_p^\top d-\kappa_p]_+^2,\\
 \text{subject to}\quad&h_y^\top d\le\epsilon_y\quad(y\in\mathcal Y_{\mathrm{keep}}),
       \qquad\|d\|_2\le\rho.\notag
\end{align}
이는 선택한 기능 방향의 증가에 벌점을 부여하고,
보존 클래스의 손실 증가도 제한하려는 국소 최적화다.
삭제 데이터로부터 제어 기준을 다시 만들거나 억제 업데이트를 되넣지 않는다.
FU의 실제 보정 갱신에 적용하며 후보 언러닝 모델을 여러 개 생성하는 탐지 과정은 아니다.
그러나 일반적인 투영·정규화 자체도 선행 기법이다. 신규성 후보는 FL에서 찾은 잠복 관계와
실제 삭제 경계가 제어 좌표를 정하고, 그 선택이 이력 없는 제어보다 나은지에 있다.

$\epsilon_y\ge0$, $\rho\ge0$이면0갱신이 제약을 만족하지만, 항상0을 내는 방어는 유용하지 않다.
진행량, 정상 성능, 망각을 함께 검증해야 하며 이 제약이 정확한 FU나 ASR 감소를 보장하지 않는다.
보존 클래스 기준 자료와 의미 보존 변환의 추가 가정, 기울기 및 최적화 비용을 실험에 포함한다.
현재 PyTorch 어댑터는 buffer 없는 모델에 한정한다. BN 상태를 무시하면 제어 방향과 실제
상태 변화가 달라질 수 있으므로 지원하지 않는 모델은 오류로 처리한다.
원래 BadFU에서 과거 표본별 장부를 얻지 못하면 이력 기반 좌표 선택도 미확정이다.
이때 새 표본 장부만을 유일한 해결책으로 가정하지 않는다.
실제 FU에서 camouflage를 제거한 보존 데이터로 계산한 클라이언트 보정 갱신은 이미
추가 관측을 제공할 수 있다. 동일한 현재 모델의 Jacobian $J_q$에서
$J_q u_j^{\mathrm{FL}}$와 $J_q\widetilde u_j^{\mathrm{FU}}$를 비교하면,
과거 혼합 갱신과 요청 이후 보정 갱신의 기능 응답 차이를 측정할 수 있다.
새 삭제 시뮬레이션이나 과거 표본 기여의 정확한 분해를 요구하는 관측은 아니다.
단, 학습 길이·노름·모델 시점의 비교 조건을 맞춰야 하며 정상적인 표본 삭제도 차이를 만든다.
이 관측에서 선택한 방향과 이력만으로 선택한 방향을 각각 소거해야 한다.
원래 BadFU의 내부 역할을 자동 복원하거나 그것의 방어 성공을 이미 보장한다는 뜻은 아니다.

\subsection{국소 최적화의 성질과 정확한 보장 범위}
투영 자체는 새로운 발명이 아니다. GEM은 과거 과제 손실 기울기에 대한 제약으로
가까운 업데이트를 구하는 선행 방법이다\cite{gem}. 우리 검증 대상은
식~\eqref{eq:action-selection}이 이력 없는 같은 제어보다 유용한 방향을 고르는가다.

\begin{proposition}[계산 부분공간]
식~\eqref{eq:controlled-step}의 최적해는
$\operatorname{span}\{v,g_p,h_y\}$에 존재한다.
\end{proposition}
\begin{proof}
임의 갱신의 해당 공간에 직교하는 성분은 모든 기능 및 손실의 선형 반응을 바꾸지 않는다.
그 성분을 제거하면 거리 목적과 노름은 증가하지 않고 모든 제약이 보존된다.
\end{proof}
구현은 SVD의 수치 rank로 최대 $1+|\mathcal P_r|+|\mathcal Y|$차원으로 줄여 SLSQP를 실행한다.
파라미터 수 $D$에 대해 기저 메모리는 $O(D(M+C))$, 기저 계산은 $O(D(M+C)^2)$다.
수치 성공 및 제약 잔차를 검사하며 solver 실패를0갱신 성공으로 대체하지 않는다.

\begin{proposition}[조건부 국소 증가 상한]
$\lambda>0$, $\kappa_p,\epsilon_y,\rho\ge0$이고 정확한 최적해를 구했다고 하자.
반지름 $\rho$에서 $q_p$의 기울기가 $L_p$-Lipschitz이면 선택된 좌표에 대해
\begin{equation}
q_p(\theta+d^*)-q_p(\theta)
\le\kappa_p+\frac{\|v\|}{\sqrt{2\lambda}}+\frac{L_p\rho^2}{2}.
\end{equation}
보존 손실이 같은 영역에서 $L_y$-smooth이면 그 증가는 $\epsilon_y+L_y\rho^2/2$ 이하이다.
\end{proposition}
\begin{proof}
0은 가능해이므로 최적 목적값은 $\|v\|^2/2$ 이하이다.
따라서 각 양의 hinge 항은 $\|v\|/\sqrt{2\lambda}$ 이하이다.
이를 smoothness 부등식에 대입한다. 보존 손실에는 선형 제약을 대입한다.
\end{proof}
알 수 없는 Lipschitz 상수를 수치 인증으로 사용하지 않으며 ReLU 네트워크에 전역 smoothness를
자동 가정하지 않는다. 이 명제는 전체 ASR, 선택되지 않은 트리거, 반복 FU 또는 망각 보장이 아니다.

\subsection{실제 한 단계의 개발 검사와 관측 상한}
MNIST/LeNet 파라미터44,426개, K20, seed522의 기존 한 궤적을 사용했다.
최근20개 전이의 FL 기록, 선택32개, 보존128개, 별도 평가512개를 사용한다.
99라운드의 실제 참여8명에서 요청자4를 제외한7명이 각2batch SGD 보정을 수행한다.
$\lambda=1$, $\kappa=\epsilon=0$, $\rho=.25\|v\|$, $M=8$로 고정했다.
시작 상태는 global98로서 이미 삭제 기여를 포함하므로 이 실험은 전체 망각이 아니다.
기본 관측은 밝기 $.9x$, 추가 진단은 저장된 실제 트리거로 만든 입력이다.
정답 트리거 진단에서 장력 $B$의 상위8좌표는 모두 공격 표적8이었다.
밝기 관측의 상위8좌표에는 표적8이 없었다. 표적 번호는 점수 계산에 주지 않았지만,
트리거 입력 자체는 oracle 정보이므로 실제 미지 트리거 방어 성능으로 보고하지 않는다.

\begin{table}[ht]\centering\small
\caption{R-039: 동일 보정 한 단계의 개발 결과. ASR 분모는 표적 제외458개. 전체 FU 성능이 아니다.}
\begin{tabular}{lrr}\toprule
조치 & 단계 후 ASR(\%) & ACC(\%)\\\midrule
단계 전 &12.45&98.83\\
무방어 &20.74&88.67\\
clipping만 &12.23&98.83\\
보존 제약만 &11.35&99.02\\
밝기 관측, 이력 없는 제어 &11.79&99.02\\
밝기 관측, 선형 영향$\times A$ &10.70&99.02\\
밝기 관측, 장력 $B\times A$ &10.70&99.02\\
oracle 트리거 관측, 장력 $B\times A$ &11.35&99.02\\\bottomrule
\end{tabular}
\end{table}
제어 실패라는 결론도, 장력 이력의 추가 이득이라는 결론도 이 표에서 도출하지 않는다.
밝기 관측의 이력 없는 제어와의 ASR 차이는5개 입력이며 독립 시드의 통계적 근거가 아니다.
clipping과 보존 제약만으로 큰 변화가 생겼고 단순 선형 기준과 제안은 같은 ASR이다.
oracle 장력의 표적 식별과 최종 ASR 감소 역시 다른 질문이다.
100단계 및 수정 전 개발 로그도 보존했으며 같은 궤적의 여러 단계를 독립 시드로 세지 않는다.

\paragraph{망각과 안전성의 목적 충돌.}
보존 데이터에 공격이 남아 있으면 일반 재학습도 백도어를 학습할 수 있다.
일반 재학습과 동일한 방어 목적을 적용한 보존 데이터 재학습을 모두 비교해야 한다.
후자와의 근접성만으로도 형식적 망각 인증이 되지는 않는다.
삭제 유래 이력으로 고른 좌표 자체에 삭제 데이터 의존성이 남을 수 있으므로,
원시 장부 폐기를 그 의존성 제거의 증명으로 쓰지 않는다.
엄격한 망각을 주장하려면 유도 상태까지 포함한 별도 삭제 분석 또는 보존 정보만의 기준 재구성이 필요하다.

\subsection{다음 실험은 연결 가정과 신규성을 분리한다}
먼저 기존 체크포인트로 다음 세 관문을 확인한 뒤 독립 학습을 확대한다.
\begin{enumerate}
\item \textbf{관측 관문:} 원시 장력, 입력 쌍 장력, 단순 선형 삭제 영향에 대해
      동일 정상 예산의 분리와 삭제 경계 순위를 비교한다. 실제 트리거 상한에서조차
      분리가 없으면 임의 변환 탐지기의 고도화를 먼저 진행하지 않는다.
\item \textbf{조치 관문:} 같은 FU 갱신에 대해 무방어, 기존 clipping/방어,
      이력 없는 식~\eqref{eq:controlled-step}, 삭제 경계만 사용하는 제어,
      FL 이력과 경계를 모두 사용하는 제어를 비교한다.
      이력의 역할을 없애거나 섞어도 성능이 같으면 과정 연결의 기여를 주장하지 않는다.
\item \textbf{효용 관문:} 후속 ASR, 보존 ACC와 희귀 클래스 ACC, 전체 요청 망각,
      연산 비용, 이후 재학습에서의 재활성화를 함께 측정한다.
      공격 없는 non-IID의 자연스러운 오분류와 공격의 추가 효과를 따로 보고한다.
\end{enumerate}
평가에서는 정답 트리거 $g$와 원래 표적이 아닌 출처 클래스 $y\ne c$를 사용한다.
시점 $t$의 예측을 $\widehat y_t(x)=\arg\max_df_d(x;\theta_t)$로 두고,
$A^{g,t}_{y\to c}=\Pr[\widehat y_t(g(x))=c\mid y]$와
$A^{0,t}_{y\to c}=\Pr[\widehat y_t(x)=c\mid y]$를 모두 측정한다.
\begin{equation}
 \Delta_{\mathrm{act}}=\frac1{C-1}\sum_{y\ne c}
 \big\{(A^{g,\mathrm{FU}}_{y\to c}-A^{0,\mathrm{FU}}_{y\to c})
       -(A^{g,\mathrm{FL}}_{y\to c}-A^{0,\mathrm{FL}}_{y\to c})\big\}
\end{equation}
를 보조 지표로 두되 원래 ASR도 함께 보고한다.
동일 분할·요청의 공격 없는 대조에서도 이를 계산한다.
이 차이 지표가 의도를 증명하는 것은 아니며, 평가용 트리거를 탐지기에 제공하지 않는다.

\section{현재 개정 알고리즘: 요청 조건부 상쇄 연산자}
\label{sec:operator-revision}
앞의 기능 좌표 기반 제어를 비교군으로 보존하고, 본 절을 현재 주 알고리즘 후보로 둔다.
기존 구현에는 이미 복원력·지연 상관·누적 장력·연속 경보·10/20라운드 창 비교가 있다.
특히 지연 상관은 기존6궤적에서 성과가 없었으므로 이름만 바꾸어 재사용하지 않는다.
개정점은 원시 업데이트의 요청 경계를 행렬로 표현하여 top-k 없이 제약 최적화에 연결하는 것이다.

\subsection{시간 추정과 상쇄 행렬}
실제 집계 가중치가 반영된 업데이트를 $x_{it}$라 하자. 요청 $R$에 대해
$d_t=\sum_{i\in R}x_{it}$, $k_t=\sum_{i\notin R}x_{it}$를 만들고,
행이 시간인 두 행렬을 $D,K$라 한다. 미참여자의 실제 집계 기여는0이지만
미관측 local gradient를0으로 관측했다고 해석하지 않는다.
요청 전에 존재하는 기록 창에서 각각
\begin{equation}
\widehat X=\arg\min_Z\left\{\|Z-X\|_F^2+
\tau\sum_{t=1}^{T-1}\frac{\|Z_{t+1}-Z_t\|^2}{s_{t+1}-s_t}\right\}
=(I+\tau L)^{-1}X
\end{equation}
를 계산한다. $L$은 시간 차분 Laplacian이다. 완만한 집계 성분을 추정한다는
모델 가정이며, 짧은 공격을 희석할 수 있으므로 무평활과 전체 누적을 별도 비교한다.
\begin{equation}
E=\|\widehat D\|_F^2+\|\widehat K\|_F^2,\qquad
C_R=-\frac{\widehat D^\top\widehat K+\widehat K^\top\widehat D}{E},
\quad P_R=\arg\min_{P\succeq0}\|P-C_R\|_F^2=[C_R]_+=G^\top G .
\end{equation}
$E=0$이면 $P_R=0$이다. 임의 방향 $z$에 대해
\begin{equation}
z^\top C_Rz=-\frac{2}{E}\sum_t
\langle z,\hat d_t\rangle\langle z,\hat k_t\rangle .
\end{equation}
따라서 반대로 작용한 성분에 양의 값이 생긴다. PSD 투영은 비음수 제어
측도를 얻는 Frobenius 최소거리 해다. 악성 확률이나 ASR의 추정량은 아니다.
전체 파라미터 수 $p$에 대해 $p\times p$ 행렬을 만들지 않고,
$[D;K]$의 최대$2T$차원 행공간에서 고유값 문제를 푼다.
창 메모리는 $O(Tp)$, 기저 계산은 대략 $O(pT^2+T^3)$다.

\begin{proposition}[균일 삭제 소거와 좌표 변환]
모든 단위가 같은 비율 $r$로 삭제되면 $P_R=0$이다.
공통 배율 변화에 $P_R$은 불변이고 직교 좌표 변환에 등변이며,
$\|P_R\|_2\le1$이다.
\end{proposition}
\begin{proof}
균일 삭제에서 $D=rX$, $K=(1-r)X$이고 평활이 선형이므로 $C_R\preceq0$이다.
공통 배율은 분자와 $E$에서 소거되고, 직교 변환은 행렬의 합동 변환을 준다.
마지막 부등식은 각 시간의 $2|ab|\le a^2+b^2$를 합하여 얻는다.
클라이언트별 배율이나 일반 신경망 재매개화에 대한 불변성은 아니다.
\end{proof}

\subsection{최소 수정 문제와 실제 검사}
$v$는 실제 FedEraser 보정 제안, $H$의 행은 독립 validation의 클래스 손실 gradient다.
기본 개정안은 벌점 크기를 고르는 대신 노출 예산을 직접 둔다.
\begin{equation}
\begin{aligned}
\delta^*=\arg\min_\delta\;&\tfrac12\|\delta-v\|^2\\
\text{s.t. }\;&\|G\delta\|\le\eta\|Gv\|,\quad H\delta\le\epsilon,\quad
\|\delta\|\le\rho,\quad v^\top\delta\ge\gamma\|v\|^2 .
\end{aligned}
\end{equation}
가능집합이 비어 있지 않으면 강볼록 목적에 의해 해는 유일하다.
해는 $v,G,H$의 행이 만드는 공간에 있으므로 축소 공간에서 풀 수 있다.
선택적 벌점 $\lambda\|G\delta\|^2/2$만 있고 제약이 없으면 해는
$(I+\lambda P_R)^{-1}v$다. 이 해석적 해와 제약이 충돌하는 사례를 동작 검사한다.

무차원 좌표 $z$의 목적 $F$는1-강볼록이다. 가능점, 비음수 승수 $\mu$,
볼록 제약 $f_i(z)\le0$ 및 stationarity 잔차 $r$에서
\begin{equation}
F(z)-F(z^*)\le\tfrac12\|r\|^2-\sum_i\mu_i f_i(z)
\end{equation}
를 얻는다. 강볼록 Lagrangian의 최소값 하한과 약한 쌍대성으로 증명된다.
구현은 수치 가능성 오차와 이 상계를 보고한다. solver 종료 메시지만으로
최적성이나 불가능성을 선언하지 않는다.

실제 클래스 손실 $L_c(\theta+\delta)-L_c(\theta)$도 검사한다.
위반 시 갱신을 축소하되 최소 진행 $\gamma$ 아래로 내리지 않는다.
동시에 만족하지 못하면 불수용을 보고하고 replay를 중단한다.
0 갱신을 망각 성공으로 처리하지 않는다. $\gamma$는 제안 방향 진행이며
망각 비율이 아니다. 반복 단계의 손실 예산 누적과 최종 효용은 별도 평가한다.

\subsection{구현과 현재 증거}
\texttt{request\_geometry.py}는 연산자·최적화·실제 검사·replay callback을,
\texttt{experiment\_request\_replay.py}는 초기 상태부터 각 군의 보존 업데이트를
다시 계산하는 실행을 제공한다. 진짜 학습 전 상태와 실제 보존 자료를 요구한다.
새 학습은 round0 및 참여 명단을 저장하고 정상 궤적도 local 기록을 남긴다.
기존100라운드 자료에 진짜 초기 상태와 대응 보존 archive가 없어
전체 공격 replay를 실행했다고 보고하지 않는다. 합성 자료의 실제 PyTorch
2라운드 통합 검사는 실행 경로 검사다.

기존 K20/seed522의99단계에서 $\eta=.1$, $\epsilon=.01$ nats,
$\gamma=.1$, $\rho=.25\|v\|$, $\tau=1$을 동일 적용했다.
\begin{table}[ht]\centering\small
\caption{R-044/045: 같은 한 모델의 단계 검사. ASR/ACC(\%). 독립 학습 시드가 아니다.}
\begin{tabular}{lccc}\toprule
요청 & 무방어 & 이력 없는 보존 제어 & 개정 제어\\\midrule
4 &20.74/88.67&11.14/99.02&11.14/99.02\\
8 &18.12/58.40&11.79/99.02&11.79/99.02\\\bottomrule
\end{tabular}
\end{table}
요청4에서 정의한 노출 제곱은 $3.024\cdot10^{-4}$에서
$3.024\cdot10^{-6}$으로 줄었고 실제 최대 클래스 손실 증가는.009813,
진행량은.24169였다. 요청8은 실제 검사에서 한 번 축소되어 진행량.12077,
최대 손실 증가.004425로 수용되었다. 선형·누적·무평활·순서 섞기의 ASR도 같았다.
즉 제약 최적화의 실행은 확인했지만 시간 이력의 추가 방어 효과는 미확인이다.

이차식은 요청과 여집합에 같은 상쇄 연산자를 줄 수 있고 정상 non-IID의 상쇄도
포함한다. 기능적 상쇄가 항상 큰 parameter 상쇄로 관측된다는 보장은 없다.
공통 구조를 적용해도 원래 BadFU의 표본 기록과 BN 어댑터가 자동으로 생기지 않는다.
GEM과 FedOSD도 제약 방향과 충돌 처리를 사용하므로
최적화 형식 자체를 신규성으로 주장하지 않는다\cite{gem,fedosd}.
기여 후보는 실제 삭제 경계로 구성한 상쇄 연산자가 동일 정보의 제어보다
ASR--효용--삭제 이행에 추가 가치를 주는가에 있다.

\section{결론}
FUBA와 BadFU를 연결하는 것은 주입과 억제의 공동 학습 및 억제의 선택적 삭제다.
우리는 장력을 이 구조의 학습 중 관측으로 두고,
삭제 경계가 어느 방향의 균형을 깨는지 계산하는 절차로 발전시킨다.
공격 구조의 공통성은 같은 정의로 표현하고, 관측 해상도의 차이는 구현과 검증의 문제로 다룬다.
기능 관측의 예비 결과와 원시 업데이트 연산자를 구분하며,
현재 개정안은 상쇄 노출·정상 손실·최소 진행의 명시적 제약을 실제 갱신에 적용한다.
수치 최적성 및 실제 손실 검사를 통과하는 것과 백도어 안전성은 별개다.
전체 FU에서의 ASR·보존 성능·망각, 원래 BadFU와 정상 대조 및 이력의 추가 가치가 다음 필수 검증이다.

\appendix
\section{이전 실험의 보존 범위}
다음 부록은 앞선 v7에서 중심이었던 언러닝 후 반응 검사다.
현재는 비교군 및 실제 효과 확인 자료이며 장력 기반 검출의 본 알고리즘이 아니다.
해당 실험의 근사 삭제 연산자는
\begin{equation}\label{eq:proxy}
 \theta(r)=\theta_0-\sum_jr_jU_j
\end{equation}
였다. 원래 BadFU의 데이터 묶음 삭제, 정확한 망각, 새 장력 점수에 의한 복구를 수행한 실험이 아니다.
이 부록에서 $p_r(x)=\operatorname{softmax}f(x;\theta(r))$이며,
공개 뷰와 요청 뷰의 선택·확인 입력을 각각 $D,H$로 표시한다.

\section{보조 비교군: 언러닝 후 soft 반응 검사}
\subsection{삭제 효과의 비교 기준}
전체 요청과 단위 수 및 비율이 같은 비교 요청 $r_b$를 만든다.
각 뷰의 입력을 선택 집합 $D$와 확인 집합 $H$로 나눈다.
선택 집합에서 기준 모델과 argmax가 바뀐 비율이 $0.3$ 이하인 비교 요청을 $B$에 남긴다.
이 조건은 정상성 증명이 아닌 비교 가능성 휴리스틱이다.
입력별 기준 예측을
\begin{equation}
 \widetilde q_c(x)=\med_{b\in B}p_{r_b,c}(x),\qquad
 q_c(x)=\frac{\widetilde q_c(x)}{\sum_d\widetilde q_d(x)}
\end{equation}
로 정한다. 분모가 0이면 비교 예측의 평균을 쓴다.
비교군 자체의 점수에는 자신을 뺀 $q^{(-b)}$를 사용한다.
현재 비교는 데이터 크기, 참여 횟수, 클래스 분포까지 일치시키지 않으며
비교군에 공격 단위가 포함될 수 있다.

\subsection{클래스 전환 이전의 확률 흐름}
집합 $X$에서 원래 클래스 $a(x)=\arg\max_c q_c(x)$를 정의한다.
후보 삭제가 기준보다 더 부여한 확률을 원래 클래스별로 모은다.
\begin{equation}\label{eq:flow}
 F_{a\to c}(r;X)=\frac1{|X|}\sum_{x\in X}
 \ind\{a(x)=a,\ a\ne c\}\pos{p_{r,c}(x)-q_c(x)}.
\end{equation}
특정 출력으로 유입되는 질량과 원래 클래스의 다양성을 함께 측정한다.
\begin{align}
 m_c&=\sum_{a\ne c}F_{a\to c},& \pi_{a,c}&=F_{a\to c}/m_c,\\
 h_c&=\frac{-\sum_{a\ne c}\pi_{a,c}\log\pi_{a,c}}{\log(C-1)},& S_c&=m_ch_c.
\end{align}
$m_c=0$일 때 $S_c=0$, $0\log0=0$으로 정의한다.
이는 최적수송의 고유한 이동량이 아니라 양의 확률 증가를 묶는 관측 통계다.
단일 원래 클래스의 변화는 억제하지만 단일 출처 공격도 놓칠 수 있다.

소거군 hard 흐름은 식~\eqref{eq:flow}의 양의 확률 증가를
$\ind\{\arg\max_d p_{r,d}(x)=c\}$로 바꾸고 나머지 절차를 동일하게 유지한다.
기존 DOSE는 기준과 표준화, 표적 선택 규칙에도 차이가 있으므로
soft 관측 자체의 효과는 이 hard 소거군과 비교한다.

\subsection{표적 선택과 확인}
분할 $A\in\{D,H\}$에서 비교군 점수 $S_{b,c}^A=S_c(p_{r_b},q^{(-b)};A)$로
\begin{align}
 \mu_c^A&=\med_b S_{b,c}^A,\\
 s_c^A&=\max\{1.4826\med_b|S_{b,c}^A-\mu_c^A|,\epsilon\},\\
 z_c^A&=(S_c^A-\mu_c^A)/s_c^A
\end{align}
를 계산한다. 선택 집합에서 $\widehat c=\arg\max_c z_c^D$를 정하고 확인 집합에서 바꾸지 않는다.
뷰 $v$의 경보는
\begin{equation}\label{eq:alarm}
 A_v=\bigwedge_{A\in\{D,H\}}
 \{z_{\widehat c}^A>\tau\ \land\ S_{\widehat c}^A-\mu_{\widehat c}^A>\delta\}
\end{equation}
이다. 실험값은 $\tau=3$, $\delta=0.02$, $\epsilon=0.005$로 고정했다.
뷰당 256개 입력을 seed 20261014로 128/128 분할하며, 입력32개 또는 비교군3개 미만이면 판정 불가다.
이진 과제는 엔트로피가 퇴화하므로 판정 불가로 처리한다.

z를 Gaussian p값으로 해석하지 않는다. 비교 요청은 독립 표본이 아니고,
확인 입력을 분리해도 학습 모델과 데이터 출처를 공유한다.
후보와 비교군의 기준 모델 수도 하나 차이 난다.
따라서 이 규칙에는 conformal 보장이나 지정된 오탐률이 자동으로 부여되지 않는다.

\subsection{판정에서 조치까지}
\begin{enumerate}
\item 기록과 누적 요청을 검증하고 같은 연산자로 후보 및 비교 삭제 모델을 만든다.
\item 공개·요청 뷰에서 기준, soft 흐름, 선택 표적, 확인 점수를 계산한다.
\item 어느 뷰든 식~\eqref{eq:alarm}을 만족하면 보류한다. 공개 뷰 평균 TV가
      $0.15$를 초과하는 보존 예산 위반도 별도로 보류한다.
\item 필요한 경우 CCD 이력, TRACE 후보, 실제 공동삭제를 진단 증거로 수집한다.
      설명 상대를 찾지 못했다는 이유로 다른 검사의 위험을 취소하지 않는다.
\item 전체 요청을 반영한 복구 모델을 만들고 보존·잔여 공격·망각을 독립 검증한 뒤 교체한다.
\end{enumerate}
후보 공격자 자동 배제는 수행하지 않는다. 현 어댑터에는 실제 복구와 독립 망각 검증이 없어
무경보여도 감사만 완료된 상태로 남는다. 복구 이후에는 기준 모델과 기여 기록을 갱신해야 한다.
단일 요청의 모든 동료를 비교하면 기준을 포함해 최대 $K+1$개의 모델 예측이 필요하며,
비교 조합 수와 연산자에 따라 비용이 증가한다.

\section{보조 soft 통계의 성질}
\begin{proposition}[범위와 단일 출처]
$0\le S_c\le1$이며 하나의 원래 클래스에서만 $c$로 질량이 들어오면 $S_c=0$이다.
\end{proposition}
\begin{proof}
확률분포 차이의 양의 부분 합은 TV로 1 이하이며 $0\le h_c\le1$이다.
단일 출처에서는 $\pi_{a,c}$가 한 점에 집중하여 엔트로피가 0이다.
\end{proof}
이 성질은 모든 정상 non-IID 변화가 허용된다는 정리가 아니다.
정상 삭제가 여러 원래 클래스에 영향을 주면 큰 점수를 만들 수 있다.

\begin{proposition}[argmax 관측의 사각지대]
3클래스 입력에서 두 원래 클래스가 같은 수로 나타나고 기준 예측이 각각
$(.65,.20,.15)$, $(.20,.65,.15)$라고 하자.
원래 최대 확률에서 $0<\eta<.25$를 빼 세 번째 클래스에 더하면
hard 점수는 0이지만 soft 점수 $S_3=\eta$다.
\end{proposition}
\begin{proof}
$.65-\eta>.15+\eta$이므로 argmax는 바뀌지 않는다.
두 원래 클래스에서 각각 $\eta/2$가 유입되어 질량 $\eta$와 정규화 엔트로피1을 얻는다.
\end{proof}
이는 추가 정보의 구성적 예이다. 실제 모든 모델의 검출력 증가나 ASR에 대한 하한이 아니다.

\begin{proposition}[공통 반응과 비식별성]
후보 예측이 모든 비교 요청의 예측과 같으면 $S_c=0$이다.
정상 삭제와 공격 삭제가 모든 관측 예측에서 같다면 이 관측에 기반한 판정기는
두 경우를 항상 구별할 수 없다.
\end{proposition}
\begin{proof}
첫 성질은 $p_r=q$를 식~\eqref{eq:flow}에 대입하면 얻는다.
두 번째 성질은 동일 관측에 대한 함수의 출력이 같다는 사실로 따른다.
\end{proof}
상쇄의 존재도 악의의 충분조건이 아니다. 공개·요청 입력이 트리거 영역을 관측하지 못하는 경우,
확률 증가가 한 원래 클래스에 국한되거나 여러 표적으로 분산되는 경우,
또는 비교군이 오염되는 경우는 주요 실패 가능성이다.

\section{이전 언러닝 후 검사 실험의 원기록}
\subsection{자료와 프로토콜}
결과 장부 R-028--R-031의 저장 모델 감사84조건을 사용했다.
FUBA는 MNIST/LeNet, K20, 100라운드, 참여율0.4의 IID seed522 한 궤적이고,
요청자는 억제자4 및 정상5,19다. BadFU 변형은 CIFAR-10/ResNet18,
K6, 15라운드 한 궤적에서 camouflage 요청자5 및 정상1,2를 검사했다.
이 변형은 원래 BadFU의 같은 클라이언트 내 표본 삭제 재현이 아니다.
정상 대조는 MNIST/LeNet, K8, 8라운드 한 궤적의 모든8명이다.

비율은 $.05,.10,.20$이며 두 프로브 구성을 사용했다.
기존 구성은 이전 감사 입력을 재사용하고, 새 구성은 다음 공개 입력256개와
이전 요청 프로브를 제외한 새 요청 입력256개를 사용했다.
새 공개 입력 일부는 이전 ACC 평가 구간과 겹치므로 완전히 봉인된 시험 집합이라 하지 않는다.
역할과 트리거 정답은 판정에 사용하지 않았다.
문턱과 분할 규칙은 실행 전에 고정했고 실패 후 조정하지 않았다.
반복 비율과 프로브는 독립 학습 시드가 아니다.

\subsection{확률 관측은 hard 관측에 없는 신호를 주는가}
\begin{table}[ht]\centering\small
\caption{검출기 경보 수: 기존 입력 / 새 입력. 분모는 프로브 구성당 요청--비율 수다.
보존 예산 위반은 경보 수에 합산하지 않는다.}
\begin{tabular}{lrrrr}\toprule
대상 & 분모 & 기존 DOSE & hard 확인 & soft 확인\\\midrule
FUBA 억제자 & 3 & 1 / 1 & 1 / 1 & 2 / 2\\
FUBA 정상 & 6 & 0 / 0 & 0 / 0 & 0 / 0\\
BadFU 변형 요청자 & 3 & 1 / 0 & 0 / 0 & 0 / 0\\
BadFU 변형 정상 & 6 & 0 / 1 & 0 / 0 & 0 / 0\\
공격 없는 정상 & 24 & 0 / 0 & 0 / 0 & 0 / 0\\\bottomrule
\end{tabular}
\end{table}

FUBA 억제자의 $.10$ 삭제에서 soft는 두 입력 구성 모두 공개 뷰의 클래스8을 선택했고
확인 z가 각각6.154, 6.050이었다. 같은 분할의 hard와 기존 DOSE는 무경보였다.
공개 TV는 $.0397,.0437$로 보존 예산보다 작았다.
별도 기존 평가 R-024에서 동일 삭제 모델의 ASR은 기준9.8\%에서49.8\%로 증가했다.
이 결과는 전체 출력 붕괴 전에 추가 관측 정보가 있음을 지지한다.
실제 입력의 모든 argmax가 불변이라는 주장은 아니며, 문턱을 동일 FPR에 맞춘 비교도 아니다.
$20\%$ 삭제는 모든 검출기가 경보를 냈고 공개 보존 예산도 위반했다.

\subsection{선택한 방향은 다른 입력에서도 유지되는가}
BadFU 변형의 $.20$ 삭제에서 기존 요청 입력의 soft 선택 z는4.162였지만
확인 z는$.676$으로 줄었다. hard도5.301에서1.358로 줄었다.
따라서 별도 확인은 모든 경우에 검출력을 높이지 않는다.
희소 신호의 표집 변동인지, 모델이나 비교 기준의 문제인지는 이 실험만으로 확정할 수 없다.

기존 DOSE도 기존 입력에서는 해당 요청에 경보를 냈지만 새 입력에서는 놓쳤다.
새 입력의 정상 요청자1, $.20$에서는 경보를 냈다.
soft의 정상 경보0만으로 우월성을 주장하면 BadFU 미검출이라는 비용을 감추게 된다.
모든 방법이 같은 학습 모델을 공유하므로 표의 차이를 모집단 오탐률이나 유의성으로 해석하지 않는다.

\subsection{상쇄 진단을 필수 단계로 둘 필요가 있는가}
이전 R-020의 실제 공동삭제 검증은 FUBA $.05,.10$에서 추가 신호가 있었지만,
새 soft는 $.05$를 놓쳤다. 반대로 새 soft의 $.10$ 경보는 상대 후보 선택 없이 발생했다.
공동삭제의 상대에는 정상 클라이언트가 섞였고 BadFU 변형에서는 상대를 찾지 못했다.
따라서 상쇄 진단은 보조 증거로 남기며 새 score의 일괄 우월성도 주장하지 않는다.

R-024에서 공동삭제는 FUBA $.10$의 ASR49.8\%를9.8\%로 되돌렸지만
이는 잠복 기준 모델의 ASR9.8\%와 같았다. 출력 복원은 백도어 제거가 아니다.
기존 정답 집합 제외 재학습 참조와 새 후보 기반 복구도 서로 다른 실험이다.

\subsection{동작 검증과 결과 보존}
입력 분할, argmax 사각지대, 단일 출처, 같은 동료 반응, 이진 과제의 판정 불가,
누적 요청 및 배포 경계를 포함한 당시 동작 검사47개가 통과했다.
장력·보정 검사까지69개였으며 실제 제어·클라이언트별 보정·관측 사각지대 검사를 추가했다.
이는 구현 검증이며 탐지 성능의 표본 수가 아니다.
결과에는 입력 인덱스, 출처 설정, 소스 해시를 남겼고,
실행 당시 비교기 두 버전은 해시가 일치하는 소스 스냅샷으로 보관했다.
두 버전의 차이는 이진 과제 분기이며 모든 실제 실험은10클래스다.

\begin{thebibliography}{99}
\bibitem{fedosd} Z. Pan et al. Federated Unlearning with Gradient Descent and Conflict Mitigation.
AAAI, 2025. \url{https://arxiv.org/abs/2412.20200}.
\bibitem{fuba} FUBA: Backdoor Federated Learning via Federated Unlearning.
\href{https://ieeexplore.ieee.org/document/11231135/}{IEEE Transactions on Artificial Intelligence}, DOI:10.1109/TAI.2025.3630110.
\bibitem{badfu} B. Lu et al. BadFU: Backdoor Federated Learning through Adversarial Machine Unlearning.
\href{https://arxiv.org/abs/2508.15541}{arXiv:2508.15541}, 2025.
\bibitem{masa} J. Xu, Z. Zhang, R. Hu. Identify Backdoored Model in Federated Learning via Individual Unlearning.
\href{https://arxiv.org/abs/2411.01040}{WACV}, 2025.
\bibitem{fldetector} Z. Zhang et al. FLDetector: Defending Federated Learning Against Model Poisoning Attacks via Detecting Malicious Clients.
\href{https://zaixizhang.github.io/ZaixiZhang_files/FLDetector.pdf}{KDD}, 2022.
\bibitem{flame} T. D. Nguyen et al. FLAME: Taming Backdoors in Federated Learning.
\href{https://www.usenix.org/system/files/sec22-nguyen.pdf}{USENIX Security}, 2022.
\bibitem{deepsight} P. Rieger et al. DeepSight: Mitigating Backdoor Attacks in Federated Learning Through Deep Model Inspection.
\href{https://www.ndss-symposium.org/wp-content/uploads/2022-156-paper.pdf}{NDSS}, 2022.
\bibitem{unlearnguard} W. Wang et al. Poisoning Attacks and Defenses to Federated Unlearning.
\href{https://arxiv.org/abs/2501.17396}{WWW Companion / arXiv:2501.17396}, 2025.
\bibitem{pcgrad} T. Yu et al. Gradient Surgery for Multi-Task Learning.
\href{https://arxiv.org/abs/2001.06782}{NeurIPS}, 2020.
\bibitem{gem} D. Lopez-Paz, M. Ranzato. Gradient Episodic Memory for Continual Learning.
\href{https://papers.nips.cc/paper/2017/hash/f87522788a2be2d171666752f97ddebb-Abstract.html}{NIPS}, 2017.
\bibitem{federaser} G. Liu, X. Ma, Y. Yang, C. Wang, J. Liu.
FedEraser: Enabling Efficient Client-Level Data Removal from Federated Learning Models.
\href{https://arxiv.org/abs/2012.13891}{IWQoS, preprint: Federated Unlearning}, 2021.
\bibitem{fedredefense} Y. Xie, M. Fang, N. Z. Gong.
FedREDefense: Defending against Model Poisoning Attacks for Federated Learning using Model Update Reconstruction Error.
\href{https://proceedings.mlr.press/v235/xie24c.html}{ICML}, 2024.
\bibitem{indicator} BackdoorIndicator: Leveraging OOD Data for Proactive Backdoor Detection in Federated Learning.
\href{https://arxiv.org/abs/2405.20862}{arXiv:2405.20862}, 2024.
\end{thebibliography}
\end{document}
